\chapter*{Notation}
\addcontentsline{toc}{chapter}{Notation}
\label{chap:notation}
\markboth{Notation}{Notation}

% One group of symbols: #1 = title, rows are "symbol & description \\".
\newenvironment{notationgroup}[1]{%
  \par\medskip\renewcommand{\arraystretch}{1.3}
  \begin{center}\textbf{#1}\end{center}
  \begin{longtable}{@{}>{\centering\arraybackslash}p{0.25\linewidth}p{0.69\linewidth}@{}}
}{%
  \end{longtable}
}

This section summarizes the notation used throughout this thesis.
It follows the conventions of Goodfellow et al.~\parencite{goodfellow2016deep} for deep learning and of Sutton and Barto~\parencite{sutton2018reinforcement} for reinforcement learning.
Vectors and matrices are written in bold upright letters and sets in calligraphic letters.
Unlike Sutton and Barto, random variables are written in lower case, like their values, and the reward $r_t$ is indexed by the time step of the action $a_t$ that produced it.
Unlike Goodfellow et al., $\sigma$ denotes any activation function and $\eta$ the learning rate.

\begin{notationgroup}{Numbers and Arrays}
$a$ & A scalar \\
$\vect{a}$ & A vector \\
$\mat{A}$ & A matrix \\
$a_i$ & Element $i$ of vector $\vect{a}$, with indexing starting at 1 \\
$A_{i,j}$ & Element $i, j$ of matrix $\mat{A}$ \\
$\vect{e}^{(i)}$ & Standard basis vector, with a 1 at position $i$ \\
\end{notationgroup}

\begin{notationgroup}{Sets}
$\mathcal{A}$ & A set \\
$\mathbb{R}$ & The set of real numbers \\
$[a, b]$ & The real interval including $a$ and $b$ \\
$|\mathcal{A}|$ & Number of elements of the set $\mathcal{A}$ \\
$\Distributions{\mathcal{A}}$ & Set of probability distributions over $\mathcal{A}$ \\
\end{notationgroup}

\begin{notationgroup}{Calculus, Probability and Functions}
$\dfrac{\partial y}{\partial x}$ & Partial derivative of $y$ with respect to $x$ \\
$\nabla_{\params} y$ & Gradient of $y$ with respect to $\params$ \\
$\mat{J}_{f}$ & Jacobian matrix of the function $f$ \\
$x \sim P$ & $x$ is sampled from the distribution $P$ \\
$\mathbb{E}[\cdot]$ & Expectation \\
$f : \mathcal{A} \to \mathcal{B}$ & The function $f$ with domain $\mathcal{A}$ and range $\mathcal{B}$ \\
$f \circ g$ & Composition of the functions $f$ and $g$ \\
$\hat{f}$ & An approximation of $f$; the hat denotes an approximation \\
$\|\vect{x}\|_p$ & $L^p$ norm of $\vect{x}$ \\
$\argmax_{a} f(a)$ & A value of $a$ at which $f(a)$ is maximal \\
$\sigma$ & Activation function, applied element-wise \\
\end{notationgroup}

\begin{notationgroup}{Deep Learning}
$\mathcal{X}, \mathcal{Y}$ & Input and output spaces \\
$\vect{x}, \vect{y}$ & An input and its target output \\
$\hat{\vect{y}}$ & Output predicted by the network \\
$\mathcal{D}$ & A dataset \\
$N$ & Number of examples in a dataset \\
$(\vect{x}^{(i)}, \vect{y}^{(i)})$ & The $i$-th example of a dataset \\
$\mathcal{L}$ & Loss function \\
$\loss{\params}{i}$ & Loss of the $i$-th example \\
$\params$ & Parameters of a model, a vector of $\mathbb{R}^d$ \\
$\theta_j$ & The $j$-th parameter \\
$\hat{f}_{\params}$ & Network with parameters $\params$ \\
$\params_k$ & Parameters after $k-1$ steps of gradient descent \\
$\eta$ & Learning rate \\
$\vect{g}$ & Gradient estimate used by stochastic gradient descent \\
$L$ & Index of the output layer of a network, the input layer being indexed 0 \\
$n_l$ & Number of neurons of layer $l$ \\
$\mat{W}^{(l)}, \vect{b}^{(l)}$ & Weight matrix and bias vector of layer $l$ \\
$\latentRepresentation[l]$ & Latent representation (activations) of layer $l$ \\
$h^{(l)}_i$ & Activation of neuron $i$ of layer $l$ \\
$\phi^{(l)}$ & Function computed by layer $l$ \\
$\Encoder, \Decoder$ & Encoder and decoder of an autoencoder \\
$\vect{z} \in \mathcal{Z}$ & Code produced by an encoder, and the space of codes \\
\end{notationgroup}

\begin{notationgroup}{Latent Space Analysis}
$c_j$ & Presence of concept $j$ \\
$\vect{v}_i$ & Direction associated with concept $i$ \\
$\alpha$ & Strength of a perturbation along a direction \\
$\Probe{l}{j}$ & Probe predicting concept $j$ from layer $l$ \\
$\lambda, \lambda_1, \lambda_2, \lambda_3$ & Weighting coefficients of a loss \\
\end{notationgroup}

\begin{notationgroup}{Reinforcement Learning}
$(\StateSet, \ObsSet, \ActionSet, \RewardSet, p)$ & A Markov decision process \\
$\StateSet, \ObsSet, \ActionSet, \RewardSet$ & Sets of states, observations, actions and rewards \\
$\StateSetInitial, \StateSetTerminal$ & Sets of initial and terminal states \\
$s, o, a, r$ & A state, an observation, an action and a reward \\
$s', o'$ & Next state and next observation \\
$\StateInitial, \StateTerminal$ & An initial state and a terminal state \\
$p(s', o, r \mid s, a)$ & Probability of the transition to $(s', o, r)$ from state $s$ and action $a$ \\
$\beta(s \mid o)$ & Probability of state $s$ given observation $o$ \\
$\pi(a \mid o)$ & Probability of taking action $a$ given observation $o$ under policy $\pi$ \\
$\pi_*$ & Optimal policy \\
$\varepsilon$ & Probability of taking a random action in an $\varepsilon$-greedy policy \\
$t$ & Discrete time step \\
$T$ & Final time step of an episode \\
$s_t, o_t, a_t, r_t$ & State, observation, action and reward at time step $t$ \\
$\tau$ & A trajectory \\
$\tau_{t:}$ & Part of the trajectory $\tau$ starting from time step $t$ \\
$\tau^{\mathrm{ter}}, \tau^{\mathrm{ep}}$ & A terminal trajectory and an episode \\
$\tau \sim \pi$ & Trajectory generated by following the policy $\pi$ \\
$\Pr_\pi(\tau \mid s, o)$ & Probability of $\tau$ under $\pi$, starting from state $s$ and observation $o$ \\
$\gamma$ & Discount factor \\
$G_t$ & Return following time step $t$ \\
$v_\pi(s, o)$ & Value of state $s$ with observation $o$ under policy $\pi$ \\
$q_\pi(s, a)$ & Value of taking action $a$ in state $s$ under policy $\pi$ \\
$q_*$ & State-action value function of the optimal policy \\
$\hat{v}(s, o, \weights)$ & Approximate state value given the weight vector $\weights$ (critic) \\
$\hat{q}(o, a, \weights)$ & Approximate state-action value given the weight vector $\weights$ \\
$\pi_{\params}$ & Policy with parameters $\params$ (actor) \\
$\pi_k, \hat{q}_k$ & Policy and value approximation at iteration $k$ \\
$A(s, a)$ & Advantage of action $a$ in state $s$ \\
\end{notationgroup}

\begin{notationgroup}{Clustering Agent Representation}
$\InitialPolicy$ & Initial policy to be explained \\
$\SurrogatePolicy{i}$ & The $i$-th surrogate policy, with parameters $\params_i$ \\
$n$ & Number of surrogate policies \\
$\SurrogateEncoder[i], \Decoder[i]$ & Encoder and decoder of the $i$-th surrogate policy \\
$\logits$ & Parameters (logits) of an action distribution \\
$\ObsData$ & Set of observations used for training and clustering \\
$\LatentSet{i}$ & Latent representations of $\ObsData$ given by $\SurrogateEncoder[i]$ \\
$\ClusteringFunction$ & Clustering function \\
$\mathcal{C}_i$ & Clustering of $\LatentSet{i}$ \\
$K$ & Number of clusters \\
$\BaselineClustering{obs}, \BaselineClustering{act}$ & Clusterings of the observations and of the action distributions \\
$\BaselineClustering[u]{untr}$ & Clustering given by the $u$-th of $U$ untrained networks \\
$\Similarity{obs}{i}, \Similarity{act}{i}, \Similarity{untr}{i}$ & Similarity of $\mathcal{C}_i$ to the three baseline clusterings \\
$M$, $G^{(e)}_{\pi}$ & Number of evaluation episodes, and return of $\pi$ in episode $e$ \\
$\RI, \ARI$ & Rand Index and Adjusted Rand Index \\
$\BR, \CU, \AS, \ACU$ & \BehaviorReconstructionName, \ClusteringUniversalityName, \AbstractionScoreName, \AdjustedClusteringUniversalityName \\
\end{notationgroup}
